\documentclass[11pt,a4paper]{article}
\input{../../sablony/pripravy/priprava.sty}
\usepackage{tikz}
\nastavPripravu{1. ročník}{Souhrnné opakování}{136}{Matematická únikovka: finále prvního ročníku}
\begin{document}
\blok{Úvod a připomenutí}
Šest stanovišť opakuje \emph{logiku, lomené výrazy, funkce, soustavy, absolutní hodnotu a kvadratické rovnice}. Každé stanoviště vydá jednu číslici. Studenti skládají šestimístný kód. Připravený list \texttt{136\_PL\_MatematickaUnikovka.pdf} rozdávej bez učitelského řešení. Pro práci ve skupině 2--3 žáků doporučeno 5 minut na stanoviště.
\blok{Společně řešené příklady}
\begin{prikladReseniBox}[title={Příklad 1}]
Každá odpověď je \emph{jedna číslice 0 až 9}. Nelze-li získat právě jednu číslici, je nutné zkontrolovat postup i podmínky. Předej studentům jejich list a nedávej jim předem správné výsledky.
\tcblower
\emph{Výsledný kód je až v učitelském klíči pod stanovišti.}
\end{prikladReseniBox}
\begin{prikladReseniBox}[title={Příklad 2}]
V množině $\{1;2;\dots;9\}$ spočítejte čísla, která jsou sudá \emph{nebo} dělitelná $3$, ale neplatí obojí současně. Číslicí je jejich počet.
\tcblower
\[
\begin{aligned}
A&=\{2;4;6;8\},\quad B=\{3;6;9\}\\
A\triangle B&=\{2;3;4;8;9\}\\
\text{číslice}&=5
\end{aligned}
\]

\end{prikladReseniBox}
\begin{prikladReseniBox}[title={Příklad 3}]
Najděte $x$ z rovnice $\dfrac{x^2-9}{x-3}=7$. Číslicí je kořen.
\tcblower
\[
\begin{aligned}
x&\ne3\\
\frac{(x-3)(x+3)}{x-3}&=7\\
x+3&=7\\
\text{číslice}&=4
\end{aligned}
\]

\end{prikladReseniBox}
\begin{prikladReseniBox}[title={Příklad 4}]
Přímka prochází body $[1;8]$ a $[3;4]$. Určete její hodnotu v bodě $x=2$. Číslicí je $f(2)$.
\tcblower
\[
\begin{aligned}
a&=\frac{4-8}{3-1}=-2\\
f(x)&=-2x+10\\
\text{číslice}&=f(2)=6
\end{aligned}
\]

\end{prikladReseniBox}
\begin{prikladReseniBox}[title={Příklad 5}]
Vyřešte $\begin{cases}x+y=9,\\2x-y=3.\end{cases}$ Číslicí je hodnota $y$.
\tcblower
\[
\begin{aligned}
3x&=12\ \Rightarrow\ x=4\\
y&=9-4=5\\
\text{číslice}&=5
\end{aligned}
\]

\end{prikladReseniBox}
\begin{prikladReseniBox}[title={Příklad 6}]
Rovnice $|x-4|=3$ má dva kořeny. Číslicí je jejich součet.
\tcblower
\[
\begin{aligned}
x-4&=\pm3\\
x_1&=1,\quad x_2=7\\
\text{číslice}&=1+7=8
\end{aligned}
\]

\end{prikladReseniBox}
\begin{prikladReseniBox}[title={Příklad 7}]
Rovnice $x^2-5x+6=0$ má dva kořeny. Číslicí je jejich \emph{součin}.
\tcblower
\[
\begin{aligned}
(x-2)(x-3)&=0\\
x_1&=2,\quad x_2=3\\
\text{číslice}&=2\cdot3=6
\end{aligned}
\]

\end{prikladReseniBox}
\blok{Samostatné procvičování}
\begin{zadaniReseniBox}[title={Samostatná práce 1}]
Zapište za sebou šest získaných číslic v pořadí stanovišť. Jaký vznikne kód?
\tcblower
\[
\begin{aligned}
\text{kód}&=\underbrace{5\ 4\ 6\ 5\ 8\ 6}_{\text{šest číslic}}
\end{aligned}
\]

\end{zadaniReseniBox}
\blok{Volitelná rezerva a výzva}
\poznamka{Výzva pro skupiny, které skončí dříve: navrhni další stanoviště se dvěma různými metodami řešení a jednoznačnou číslicí jako výsledkem.}
\blok{Domácí úkol}
\begin{zadaniReseniBox}[title={Domácí úkol}]
\begin{enumerate}[label=\alph*),itemsep=2pt]
\item Kolik čísel z $\{1;2;\dots;12\}$ je sudých, ale není dělitelných čtyřmi?
\item Vyřešte $\dfrac{x^2-16}{x-4}=9$.
\item Přímka prochází body $[1;3]$ a $[4;9]$. Určete $f(0)$.
\item Vyřešte soustavu $x+y=11,\ x-y=3$ a napište hodnotu $y$.
\item Určete součet kořenů $|x-3|=2$.
\item Určete součet kořenů rovnice $x^2-8x+15=0$.
\end{enumerate}
\tcblower
\begin{enumerate}[label=\alph*),itemsep=2pt]
\item
\[
\begin{aligned}
K&=\{2;6;10\}\\
\text{počet}&=3
\end{aligned}
\]
\item
\[
\begin{aligned}
x&\ne4\\
x+4&=9\\
K&=\{5\}
\end{aligned}
\]
\item
\[
\begin{aligned}
a&=\frac{9-3}{4-1}=2\\
f(x)&=2x+1\\
f(0)&=1
\end{aligned}
\]
\item
\[
\begin{aligned}
2x&=14\ \Rightarrow\ x=7\\
y&=4\\
K&=\{[7;4]\}
\end{aligned}
\]
\item
\[
\begin{aligned}
x&=1\ \text{nebo}\ 5\\
x_1+x_2&=6
\end{aligned}
\]
\item
\[
\begin{aligned}
(x-3)(x-5)&=0\\
x_1+x_2&=8
\end{aligned}
\]
\end{enumerate}
\end{zadaniReseniBox}
\poznamka{Čas: 5 minut zadání pravidel, šest stanovišť po 5 minutách, 8 minut závěrečná společná kontrola, 2 minuty reflexe. Záměrně se nehodnotí známkou. Hlavní cíl: rozpoznat vhodnou metodu a spolupracovat.}
\end{document}
